\section{Particle-Mesh Gravitational Solver}

The Particle-Mesh (PM) method solves gravitational N-body dynamics
using a grid-based approach. Instead of computing pairwise forces
($O(N^2)$), it deposits particle masses onto a regular grid, solves
Poisson's equation via FFT, and interpolates forces back. This yields
$O(N + M \log M)$ complexity where $M = N_{\text{grid}}^3$, enabling
real-time simulation of millions of particles.

\section{Poisson's Equation}

Gravity is governed by the Poisson equation relating the gravitational
potential $\Phi$ to the mass density $\rho$:

\begin{equation}
\nabla^2 \Phi = 4\pi G \rho
\end{equation}

In Fourier space this becomes algebraic:

\begin{equation}
\hat{\Phi}(k) = -4\pi G \hat{\rho}(k) / k^2
\end{equation}

The force is then $F = -\nabla \Phi$, computed via finite differences
on the grid.

\section{Cloud-in-Cell Interpolation}

CIC (Cloud-in-Cell) distributes each particle's mass to the 8
surrounding grid cells using trilinear weights. For a particle at
position $x$ in cell $(i,j,k)$ with fractional offset $(w_x, w_y, w_z)$:

\begin{equation}
\rho_{i,j,k} += m \cdot (1 - w_x)(1 - w_y)(1 - w_z)
\end{equation}

The force interpolation uses the same weights in reverse (the adjoint
operation), which guarantees exact momentum conservation by Newton's
third law.

\section{Discrete Green's Function}

The continuum Green's function $G(k) \propto 1/k^2$ suffers from
aliasing on a discrete grid. We use the discrete form:

\begin{equation}
k_{\text{eff}}^2 = (2/h)^2 [\sin^2(\pi n_x/N) + \sin^2(\pi n_y/N) + \sin^2(\pi n_z/N)]
\end{equation}

This matches the finite-difference Laplacian exactly, eliminating
grid-scale artifacts and ensuring $\nabla^2 \Phi = 4\pi G \rho$ holds
to machine precision on the grid.

\section{Leapfrog KDK Integration}

The kick-drift-kick leapfrog scheme is symplectic, preserving the
Hamiltonian structure and giving excellent long-term energy conservation:

\begin{equation}
v_{n+1/2} = v_n + a_n \cdot \Delta t / 2
\end{equation}

\begin{equation}
x_{n+1} = x_n + v_{n+1/2} \cdot \Delta t
\end{equation}

\begin{equation}
v_{n+1} = v_{n+1/2} + a_{n+1} \cdot \Delta t / 2
\end{equation}

Periodic boundary conditions wrap particle positions to $[0, L)$.

\section{PM Algorithm Pipeline}

Each timestep executes the following stages:

\begin{itemize}
\item CIC deposit: particles $\to$ density grid $\rho(x)$
\item FFT forward: $\rho(x) \to \hat{\rho}(k)$
\item Green multiply: $\hat{\Phi}(k) = -4\pi G \hat{\rho}(k) / k_{\text{eff}}^2$
\item FFT inverse: $\hat{\Phi}(k) \to \Phi(x)$
\item Gradient: $F = -\nabla \Phi$ via central differences
\item CIC interpolation: grid forces $\to$ particle accelerations
\item Leapfrog KDK: kick-drift-kick time integration
\end{itemize}

Only one forward and one inverse FFT are needed per step.

\section{Simulation Parameters}

\param{G}{Gravitational constant (coupling strength)}
\param{N_{\text{grid}}}{Grid resolution (128, 256, or 512 per side)}
\param{L}{Box size (periodic domain length)}
\param{\Delta t}{Integration timestep}
\param{N}{Number of particles}

\section{Conservation Laws}

The CIC adjoint interpolation ensures exact conservation of linear
momentum. Total energy $E = T + V$ is conserved to high precision
by the symplectic integrator, with bounded oscillations rather than
secular drift.

\section{Cosmological Initial Conditions}

The cosmic web is the gravitationally amplified image of the primordial
density field, so realistic structure formation requires a physically
motivated initial field rather than random noise. The \textbf{Cosmological}
preset builds a Gaussian random field with a CDM-like power spectrum and
displaces particles using the Zel'dovich (first-order Lagrangian) approximation.

A white-noise field on the grid is transformed to Fourier space and shaped
into a density contrast with the matter power spectrum:

\begin{equation}
\hat{\delta}(\mathbf{k}) = \sqrt{P(k)}\,\hat{w}(\mathbf{k}), \qquad
P(k) \propto k^{n_s}\,T^2(k)
\end{equation}

where $n_s$ is the primordial spectral index and $T(k)$ a transfer function
that turns the spectrum over near the equality scale $k_\text{eq}$, suppressing
small-scale power. Each particle, initially on a uniform Lagrangian lattice at
position $\mathbf{q}$, is displaced by the Zel'dovich field:

\begin{equation}
\hat{\boldsymbol{\Psi}}(\mathbf{k}) = i\,\frac{\mathbf{k}}{k^2}\,T(k)\,\hat{\delta}(\mathbf{k}),
\qquad
\mathbf{x} = \mathbf{q} + D(a)\,\boldsymbol{\Psi}(\mathbf{q})
\end{equation}

Growing-mode peculiar velocities are seeded consistently with linear theory,
so the field already carries the coherent infall that grows the web:

\begin{equation}
\mathbf{u} = a\,H(a)\,f\,D(a)\,\boldsymbol{\Psi}(\mathbf{q}), \qquad
f \approx \Omega_m(a)^{0.55}
\end{equation}

The displacement amplitude $D(a)$ is normalised so the RMS displacement is a
fixed fraction of the mean inter-particle spacing, keeping the start in the
linear regime (displacements below the particle spacing).

\section{References}

\begin{itemize}
\item Hockney, R. W. \& Eastwood, J. W. (1988). Computer Simulation Using Particles.
\item Efstathiou, G. et al. (1985). ApJS 57, 241--260.
\item Klypin, A. \& Holtzman, J. (1997). arXiv:astro-ph/9712217.
\item Zel'dovich, Ya. B. (1970). A\&A 5, 84--89.
\item Bardeen, J. M. et al. (1986). ApJ 304, 15--61.
\item Springel, V. (2005). MNRAS 364, 1105--1134.
\end{itemize}
